% GENERATED by docs/tools/build-latex.js from docs/REPORT.md. Do not edit: change the Markdown and rebuild.
\chapter{Conclusions and Future Work}\label{ch:conclusion}

\section{Limitations and future work}

\begin{itemize}
\item The per-minute rate limit is per instance; a shared limiter would be needed at large scale.
\item Firebase App Check, email verification and sign-in with Google are not enabled.
\item The free Render tier sleeps when idle, so the first request after a pause is slow; Cloud Run with \texttt{-\allowbreak{}-\allowbreak{}min-\allowbreak{}instances 1} avoids this.
\item The Windows installer is not code-signed.
\item Model latency is outside our control; a lighter model or a nearer region helps.
\item Future: shared rate limiting, an admin dashboard for usage, more providers through the existing registry, Windows' built-in OCR engine.
\end{itemize}

\section{Conclusion}

IntelliPointer AI reduces \textquotedbl{}see something, get an answer\textquotedbl{} to one shortcut, one drag and one click, and shows the answer in a floating window without switching applications. The cloud side demonstrates a secure design built from managed services: Firebase Authentication for identity, Firestore transactions for a quota that cannot be exceeded even under concurrent load, a stateless HTTPS API on Render or Cloud Run that scales out, and secrets that never leave the server. Security is built in at every layer and each control is backed by an automated test, including integration tests against the real Firebase emulators. Privacy is a design rule, not an option: nothing a user captures or asks is stored or logged. Performance work brought the shortcut-to-crop time to roughly a quarter to a half second and cut the desktop footprint from about 397 MB to about 183 MB.
